% =============================================================================
% Guía Completa de OpenSCAD: Del Cubo al Robot Cuadrúpedo
% USS SpiderBot — Taller de Programación I, Universidad San Sebastián
% Compilar con: xelatex guia_openscad.tex
% =============================================================================

\documentclass[12pt, oneside, a4paper]{report}

% ─── Paquetes ────────────────────────────────────────────────────────────────
\usepackage[utf8]{inputenc}
\usepackage[spanish]{babel}
\usepackage{fontspec}
\usepackage{geometry}
\geometry{margin=2.5cm, top=3cm, bottom=3cm}
\usepackage{setspace}
\setstretch{1.15}
\usepackage{xcolor}
\usepackage{listings}
\usepackage{tcolorbox}
\usepackage{hyperref}
\usepackage{graphicx}
\usepackage{longtable}
\usepackage{booktabs}
\usepackage{fancyhdr}
\usepackage{titlesec}
\usepackage{parskip}

% ─── Configuración de hyperref ───────────────────────────────────────────────
\hypersetup{
    colorlinks=true,
    linkcolor=blue!70!black,
    urlcolor=blue!70!black,
    citecolor=blue!70!black,
    pdftitle={Guía Completa de OpenSCAD: Del Cubo al Robot Cuadrúpedo},
    pdfauthor={USS SpiderBot Project},
}

% ─── Configuración de listings para código OpenSCAD ──────────────────────────
\definecolor{codegreen}{rgb}{0,0.6,0}
\definecolor{codegray}{rgb}{0.5,0.5,0.5}
\definecolor{codepurple}{rgb}{0.58,0,0.82}
\definecolor{codeblue}{rgb}{0,0,0.8}
\definecolor{backcolour}{rgb}{0.95,0.95,0.92}
\definecolor{commentcolor}{rgb}{0.25,0.5,0.35}

\lstdefinestyle{openscadstyle}{
    language={},
    backgroundcolor=\color{backcolour},
    basicstyle=\small\ttfamily,
    commentstyle=\color{commentcolor}\itshape,
    keywordstyle=\color{codeblue}\bfseries,
    stringstyle=\color{codepurple},
    numberstyle=\tiny\color{codegray},
    breakatwhitespace=false,
    breaklines=true,
    captionpos=b,
    keepspaces=true,
    numbers=left,
    numbersep=5pt,
    showspaces=false,
    showstringspaces=false,
    showtabs=false,
    tabsize=2,
    frame=single,
    rulecolor=\color{lightgray},
    morekeywords={
        module, function, use, include, for, intersection_for,
        if, else, let, assert, echo, import, children,
        cube, sphere, cylinder, polyhedron,
        translate, rotate, scale, mirror, resize, color, multmatrix,
        union, difference, intersection, hull, minkowski,
        linear_extrude, rotate_extrude, text, polygon, offset, projection,
        render, surface, search, concat, str, len,
        true, false, undef, PI,
        center, r, d, h, s1, s2, r1, r2,
        points, faces, paths, convexity,
        $fn, $fa, $fs, $t, $preview, $vpr, $vpt, $vpd
    },
    sensitive=true,
    morecomment=[l]{//},
    morecomment=[s]{/*}{*/},
}

\lstnewenvironment{openscad}{\lstset{style=openscadstyle}}{}

\newtcolorbox{openscadbox}{
    colback=backcolour,
    colframe=lightgray,
    arc=2mm,
    boxrule=0.5pt,
    left=6pt,
    right=6pt,
    top=4pt,
    bottom=4pt,
}

% ─── Configuración de encabezados ────────────────────────────────────────────
\pagestyle{fancy}
\fancyhf{}
\fancyhead[LE,RO]{\thepage}
\fancyhead[RE]{\leftmark}
\fancyhead[LO]{\rightmark}
\renewcommand{\headrulewidth}{0.4pt}

% ─── Formato de títulos ─────────────────────────────────────────────────────
\titleformat{\chapter}[display]
    {\normalfont\huge\bfseries\color{blue!70!black}}
    {\chaptertitlename\ \thechapter}{20pt}{\Huge}
\titleformat{\section}
    {\normalfont\Large\bfseries\color{blue!60!black}}
    {\thesection}{1em}{}
\titleformat{\subsection}
    {\normalfont\large\bfseries\color{blue!50!black}}
    {\thesubsection}{1em}{}

% ─── Información del documento ───────────────────────────────────────────────
\title{
    \vspace{3cm}
    {\Huge\bfseries Guía Completa de OpenSCAD\\[0.5cm]
    \large Del Cubo al Robot Cuadrúpedo}\\[1.5cm]
    {\large Basada en el diseño del \textbf{USS SpiderBot}}\\[0.5cm]
    {\large Robot cuadrúpedo de 8 grados de libertad}\\[2cm]
}
\author{
    Taller de Programación I\\
    Universidad San Sebastián\\
    \texttt{https://github.com/moises-inc/USS-SPIDERBOT-solemne-3}
}
\date{2026}

\begin{document}

% ─── Portada ─────────────────────────────────────────────────────────────────
\maketitle
\thispagestyle{empty}
\newpage

% ─── Índice ─────────────────────────────────────────────────────────────────
\tableofcontents
\newpage

% ═════════════════════════════════════════════════════════════════════════════
% INTRODUCCIÓN
% ═════════════════════════════════════════════════════════════════════════════
\chapter*{Introducción}
\addcontentsline{toc}{chapter}{Introducción}

\textbf{OpenSCAD} es un modelador CAD 3D basado exclusivamente en código (CSG --- Constructive Solid Geometry). A diferencia de herramientas como Fusion 360 o TinkerCAD, aquí no arrastras mouse: escribes instrucciones y el programa las evalúa para generar geometría.

\section*{¿Por qué OpenSCAD para robótica?}
\begin{itemize}
    \item \textbf{Totalmente paramétrico}: cambias un número y todo el modelo se actualiza.
    \item \textbf{Versionable}: los archivos \texttt{.scad} son texto plano, ideales para Git.
    \item \textbf{Preciso}: las operaciones booleanas son exactas, no hay errores de manifold.
    \item \textbf{Gratuito y open-source}.
    \item \textbf{CLI headless}: puedes exportar STLs desde scripts.
\end{itemize}

\section*{Instalación}
\begin{center}
\begin{tabular}{ll}
\toprule
Sistema & Comando / Enlace \\
\midrule
Windows & Descargar installer desde \url{https://openscad.org} \\
macOS   & \texttt{brew install openscad} o descargar .dmg \\
Linux   & \texttt{sudo apt install openscad} (Debian/Ubuntu) \\
        & \texttt{sudo pacman -S openscad} (Arch) \\
\bottomrule
\end{tabular}
\end{center}

\section*{Interfaz}
Al abrir OpenSCAD ves tres paneles:
\begin{enumerate}
    \item \textbf{Editor de código} (izquierda) --- aquí escribes.
    \item \textbf{Vista 3D} (derecha) --- aquí ves el resultado.
    \item \textbf{Consola} (abajo) --- muestra errores, advertencias y outputs.
\end{enumerate}

\section*{Flujo de trabajo}
\begin{enumerate}
    \item Escribes código en el editor.
    \item Presionas \textbf{F5} --- \textit{Preview} (vista rápida, usa OpenCSG).
    \item Presionas \textbf{F6} --- \textit{Render} (cálculo CGAL completo, geometría exacta).
    \item Archivo > Exportar > Exportar como STL (o \textbf{F7}) --- obtienes el archivo para impresión 3D.
\end{enumerate}

% ═════════════════════════════════════════════════════════════════════════════
% CAPÍTULO 1
% ═════════════════════════════════════════════════════════════════════════════
\chapter{Fundamentos Geométricos}

\section{Formas primitivas 3D}
OpenSCAD tiene cuatro formas sólidas básicas:

\begin{openscad}
// Cubo: cube([ancho, fondo, alto], center);
cube([10, 20, 30]);                    // Esquina en el origen
cube([10, 20, 30], center=true);       // Centrado en el origen

// Esfera: sphere(r);
sphere(r = 10);                        // Radio 10mm

// Cilindro: cylinder(r, h);
cylinder(r = 5, h = 20);               // Radio 5mm, altura 20mm

// Cilindro con radios distintos (cono truncado):
cylinder(r1 = 5, r2 = 3, h = 10);
\end{openscad}

\section{Sistema de coordenadas}
OpenSCAD usa coordenadas cartesianas diestras:
\begin{itemize}
    \item \textbf{Eje X}: horizontal (derecha = positivo)
    \item \textbf{Eje Y}: profundidad (atrás = positivo)
    \item \textbf{Eje Z}: vertical (arriba = positivo)
\end{itemize}

\section{El parámetro \texttt{center}}
Por defecto, \texttt{cube()}, \texttt{cylinder()} y \texttt{sphere()} se posicionan con una esquina o base en el origen.

\begin{openscad}
// Sin center: el cubo crece desde (0,0,0) hacia (10,20,30)
cube([10, 20, 30]);

// Con center=true: el cubo se centra en (0,0,0)
cube([10, 20, 30], center=true);
\end{openscad}

\section{Resolución de curvas: \texttt{\$fn}, \texttt{\$fa}, \texttt{\$fs}}
Las formas curvas se aproximan con polígonos:

\begin{openscad}
// $fn controla el número de fragmentos (lados)
cylinder(r = 10, h = 20, $fn = 6);    // Hexágono — baja calidad
cylinder(r = 10, h = 20, $fn = 50);   // Calidad media
cylinder(r = 10, h = 20, $fn = 100);  // Muy suave — render lento

// También se puede definir globalmente
$fn = 50;  // Todas las curvas usarán 50 fragmentos
\end{openscad}

\Regla práctica: \texttt{\$fn = 30} para preview rápido, \texttt{\$fn = 50--80} para exportación a STL.

\section{Caso Real: Primitivas en el SpiderBot}
En \texttt{cad/cuerpo\_cuadruepdo.scad}, la placa base se define con \texttt{cylinder()}:

\begin{openscad}
// Archivo: cad/cuerpo_cuadruepdo.scad
$fn = 50;

cuerpo_r = 65.0;    // Radio del chasis circular
cuerpo_t = 3.0;     // Grosor de las placas

module placa_base_cuadruepodo() {
    difference() {
        // Placa inferior sólida: cilindro de 65mm de radio, 3mm de grosor
        cylinder(r = cuerpo_r, h = cuerpo_t, center=true);
        // Agujeros para los 4 pilares espaciadores M3
        for (a = [0, 90, 180, 270]) {
            rotate([0, 0, a])
                translate([spacer_r, 0, 0])
                    cylinder(r=1.6, h=cuerpo_t+2, center=true);
        }
    }
}
\end{openscad}

% ═════════════════════════════════════════════════════════════════════════════
% CAPÍTULO 2
% ═════════════════════════════════════════════════════════════════════════════
\chapter{Transformaciones Geométricas}

Las transformaciones mueven, rotan y escalan objetos. \textbf{No crean geometría nueva}, sino que modifican la existente.

\section{\texttt{translate([x, y, z])}}
Desplaza el objeto hijo en el espacio:
\begin{openscad}
translate([20, 10, 5])
    cube([10, 10, 10]);
\end{openscad}

\section{\texttt{rotate([x, y, z])}}
Rota alrededor del origen en grados:
\begin{openscad}
// Rota 45° alrededor de Z
rotate([0, 0, 45])
    cube([20, 5, 5]);
// Rota 90° alrededor de X (pone el cubo vertical)
rotate([90, 0, 0])
    cube([20, 5, 5]);
\end{openscad}

\section{\texttt{scale([x, y, z])}}
Escala en cada eje:
\begin{openscad}
// No uniforme: estira en X
scale([2, 1, 1]) sphere(r = 10);
// Uniforme: agranda todo al doble
scale(2) sphere(r = 10);
\end{openscad}

\section{\texttt{mirror([x, y, z])} --- Reflejo especular}
Refleja el objeto como si lo miraras a través de un plano:
\begin{openscad}
// Reflejar en X (plano YZ)
mirror([1, 0, 0]) cube([10, 20, 5]);
// Reflejar en diagonal
mirror([1, 1, 0]) cube([10, 20, 5]);
\end{openscad}

\section{\texttt{color("nombre")} --- Color para preview}
Asigna color a los hijos (solo Preview F5):
\begin{openscad}
color("RoyalBlue") cube([10, 10, 10]);
color([1.0, 0.5, 0]) sphere(r = 5);      // RGB [0-1]
color("LightGray", 0.5) cube([20, 5, 5]); // Semitransparente
\end{openscad}

\section{\texttt{resize([x, y, z])} --- Redimensionar}
Fuerza dimensiones absolutas, a diferencia de \texttt{scale()}:
\begin{openscad}
// Esfera de radio 10 → elipsoide de 30x60x10
resize([30, 60, 10]) sphere(r = 10);
\end{openscad}

\section{Composición: se lee de adentro hacia afuera}
\begin{openscad}
translate([30, 0, 0])
    rotate([0, 0, 45])
        cube([10, 10, 10]);
\end{openscad}

\section{Caso Real: Rotaciones y traslaciones}
En \texttt{cad/cuerpo\_cuadruepdo.scad}:
\begin{openscad}
for (a = [45, 135, 225, 315]) {
    rotate([0, 0, a])
        translate([cuerpo_r - 7, 0, cuerpo_t/2])
            rotate([0, 0, -90])
                soporte_servo_cadera();
}
\end{openscad}

\section{Caso Real 2: Código de colores}
En \texttt{cad/ensamble\_cuadruepodo.scad}:
\begin{openscad}
module dummy_servo() {
    color("RoyalBlue") {
        cube([12.5, 23.0, 22.5], center=true);
    }
    color("White") {
        translate([0, 5.5, 14])
            cylinder(r=6, h=1.5, center=true);
    }
}
\end{openscad}

% ═════════════════════════════════════════════════════════════════════════════
% CAPÍTULO 3
% ═════════════════════════════════════════════════════════════════════════════
\chapter{Operaciones Booleanas (El Corazón del Diseño)}

Las operaciones booleanas combinan sólidos para crear formas complejas.

\section{\texttt{union()}}
Fusiona todos los hijos en un solo sólido:
\begin{openscad}
union() {
    cube([10, 10, 10]);
    translate([5, 5, 0])
        cylinder(r = 5, h = 10);
}
\end{openscad}

\section{\texttt{difference()}}
Resta los hijos subsiguientes del primero:
\begin{openscad}
difference() {
    cube([20, 20, 10], center=true);
    translate([0, 0, -1])
        cylinder(r = 3, h = 12);
}
\end{openscad}

\section{\texttt{intersection()}}
Devuelve solo el volumen compartido:
\begin{openscad}
intersection() {
    sphere(r = 10);
    cube([10, 10, 10], center=true);
}
\end{openscad}

\section{Anidamiento}
Las booleanas se anidan sin límite:
\begin{openscad}
difference() {
    union() {
        cylinder(r = 20, h = 5);
        translate([15, 0, 0])
            cube([10, 10, 5], center=true);
    }
    cylinder(r = 3, h = 10, center=true);
}
\end{openscad}

\section{Caso Real: El soporte de servo}
El módulo \texttt{soporte\_servo\_cadera()} en \texttt{cad/cuerpo\_cuadruepdo.scad}:

\begin{openscad}
module soporte_servo_cadera() {
    difference() {
        // SÓLIDO BASE
        translate([0, 0, 7.5])
            cube([36, 28, 15], center=true);
        // 1. Cavidad dual-depth para servo SG90
        translate([0, 0, 2.9])
            cube([30.0, 24.2, 4.2], center=true);
        translate([0, 0, 9.5])
            cube([25.0, 24.2, 9.0], center=true);
        // 2. Ranura para bridas
        translate([0, 5.5, 7.5])
            cube([34.5, 3.2, 13.5], center=true);
        // 3. Agujeros pasantes para tornillos M2
        translate([-14.25, 0, 7.5])
            rotate([90, 0, 0])
                cylinder(r=1.0, h=35, center=true);
        translate([14.25, 0, 7.5])
            rotate([90, 0, 0])
                cylinder(r=1.0, h=35, center=true);
        // 4. Abertura para corona de salida
        translate([5.5, 12.0, 7.5])
            rotate([90, 0, 0])
                cylinder(r=6.2, h=15, center=true);
        // 5. Rebaje para rotación del fémur
        translate([6.5, 12.0, 7.5])
            cube([23, 5, 16], center=true);
    }
}
\end{openscad}

\textbf{Dual-depth pocket:} la cavidad tiene dos anchos (30mm abajo para cables, 25mm arriba para rosca de tornillos M2).

% ═════════════════════════════════════════════════════════════════════════════
% CAPÍTULO 4
% ═════════════════════════════════════════════════════════════════════════════
\chapter{Variables y Parámetros}

\section{Declaración y asignación}
En OpenSCAD las variables se asignan con \texttt{=} y son \textbf{inmutables}:
\begin{openscad}
ancho = 20; profundidad = 30; alto = 10;
cube([ancho, profundidad, alto]);
\end{openscad}

\section{Variables paramétricas}
\begin{openscad}
longitud = 50; ancho = 20; altura = 10;
pared = 2.0; tornillo_r = 1.5;

difference() {
    cube([longitud, ancho, altura]);
    translate([pared, pared, -1])
        cube([longitud - 2*pared, ancho - 2*pared, altura + 2]);
}
\end{openscad}

\section{\texttt{let()} --- Variables locales}
\begin{openscad}
let (ancho = 30, largo = ancho * 1.5) {
    cube([ancho, largo, 5]);
}
\end{openscad}

\section{\texttt{assert()} --- Validación}
\begin{openscad}
module soporte(alto, ancho) {
    assert(alto > 0, "El alto debe ser positivo");
    assert(ancho > 0, "El ancho debe ser positivo");
    cube([ancho, ancho, alto]);
}
\end{openscad}

\section{\texttt{function} --- Funciones definidas por el usuario}
\begin{openscad}
function area_circulo(radio) = PI * radio * radio;
function volumen_cilindro(radio, alto) = area_circulo(radio) * alto;
echo(volumen_cilindro(10, 20));  // 6283.19
\end{openscad}

\section{Funciones de string y lista}
\begin{openscad}
// str() — convertir a string
etiqueta = str("Pieza-", 3, ".stl");  // "Pieza-3.stl"
// concat() — unir listas
c = concat([1, 2, 3], [4, 5]);       // [1, 2, 3, 4, 5]
// len() — longitud
largo = len([10, 20, 30]);           // 3
\end{openscad}

\section{Operador ternario}
\begin{openscad}
usar_agujeros = true;
resolucion = (usar_agujeros ? 50 : 20);
\end{openscad}

\section{Constantes especiales}
\begin{openscad}
radio = 10;
circunferencia = 2 * PI * radio;  // 62.8318
echo(circunferencia);
\end{openscad}

\section{Caso Real: Variables paramétricas del SpiderBot}
En \texttt{cad/eslabon\_pata\_cuadrupedo.scad}:
\begin{openscad}
link_len = 55.0; // Longitud entre cadera y rodilla
link_h = 25.0;   // Altura extremo de rodilla
link_t = 16.0;   // Grosor bloque servo rodilla
servo_w = 23.0;  servo_t = 12.5;
servo_h = 22.5;  ear_dist = 28.5;
\end{openscad}

Si cambias \texttt{link\_len = 55} a \texttt{70}, el bolsillo del servo y la abertura del engranaje se desplazan automáticamente.

% ═════════════════════════════════════════════════════════════════════════════
% CAPÍTULO 5
% ═════════════════════════════════════════════════════════════════════════════
\chapter{Módulos --- Funciones del Diseño Paramétrico}

\section{Sintaxis básica}
\begin{openscad}
module nombre(param1, param2 = valor_defecto) {
    // Geometría aquí
}
\end{openscad}

\section{Ejemplo: tornillo paramétrico}
\begin{openscad}
module tornillo(longitud, diametro = 4) {
    radio = diametro / 2;
    cylinder(r = radio, h = longitud);
    translate([0, 0, longitud])
        cylinder(r = radio * 1.5, h = 3);
}
tornillo(20);                   // L=20, D=4
tornillo(30, 6);                // L=30, D=6
tornillo(diametro=5, longitud=15);  // Argumentos nombrados
\end{openscad}

\section{\texttt{children()} --- Módulos que envuelven geometría}
\begin{openscad}
module envolvente(radio_borde = 2) {
    minkowski() {
        children();  // Toma lo que se pase entre llaves
        sphere(r = radio_borde);
    }
}
envolvente(3) {
    cube([20, 20, 5], center=true);
}
\end{openscad}

\section{DRY en CAD}
\begin{openscad}
// MAL: código duplicado
cube([10, 20, 5]);
translate([15, 0, 0]) cube([10, 20, 5]);

// BIEN: con módulo
module bloque() { cube([10, 20, 5]); }
bloque();
translate([15, 0, 0]) bloque();
\end{openscad}

\section{Caso Real: Los módulos del SpiderBot}
\begin{openscad}
module soporte_servo_cadera() { ... }   // bracket de cadera
module placa_base_cuadruepodo() { ... } // chasis inferior
module placa_superior() { ... }         // chasis superior con cunas
module eslabon_completo() { ... }       // fémur completo
module tibia() { ... }                  // pierna inferior

// Composición en ensamble_cuadruepodo.scad:
module pata_articulada(angulo_cadera=0, angulo_rodilla=0) {
    dummy_servo();
    rotate([0, angulo_cadera, 0])
        eslabon_completo();
    translate([55, 2.75, 0])
        rotate([0, angulo_rodilla, 0])
            tibia();
}
\end{openscad}

% ═════════════════════════════════════════════════════════════════════════════
% CAPÍTULO 6
% ═════════════════════════════════════════════════════════════════════════════
\chapter{Iteración y Simetría --- El Bucle \texttt{for}}

\section{Sintaxis}
\begin{openscad}
for (i = [0 : 1 : 5]) {             // Rango [inicio:incremento:fin]
    translate([i * 10, 0, 0])
        cube([5, 5, 5]);
}
for (i = [10, 20, 30, 50]) { ... }  // Lista explícita
\end{openscad}

\section{Distribución radial}
\begin{openscad}
for (angulo = [0 : 60 : 300]) {
    rotate([0, 0, angulo])
        translate([40, 0, 0])
            cylinder(r = 2, h = 10, center=true);
}
\end{openscad}

\section{\texttt{intersection\_for()}}
\begin{openscad}
intersection_for(n = [0 : 5]) {
    rotate([0, 0, n * 60])
        translate([5, 0, 0])
            cylinder(r = 3, h = 10);
}
\end{openscad}

\section{List Comprehensions}
\begin{openscad}
cuadrados = [for (i = [0:10]) i * i];
pares = [for (i = [0:20]) if (i % 2 == 0) i];
\end{openscad}

\section{Caso Real: Las 4 patas del SpiderBot}
\begin{openscad}
// Distribuye 4 soportes en 45°, 135°, 225°, 315°
for (a = [45, 135, 225, 315]) {
    rotate([0, 0, a])
        translate([cuerpo_r - 7, 0, cuerpo_t/2])
            rotate([0, 0, -90])
                soporte_servo_cadera();
}

// Agujeros de pilares (0°, 90°, 180°, 270°)
for (a = [0, 90, 180, 270]) {
    rotate([0, 0, a])
        translate([spacer_r, 0, 0])
            cylinder(r=1.6, h=cuerpo_t+2, center=true);
}

// For anidado para pasacables
for (x = [-40, 40]) {
    for (y = [-39, 39]) {
        translate([x, y, 0])
            cylinder(r=7.0, h=30, center=true);
    }
}
\end{openscad}

% ═════════════════════════════════════════════════════════════════════════════
% CAPÍTULO 7
% ═════════════════════════════════════════════════════════════════════════════
\chapter{Operaciones Avanzadas}

\section{\texttt{hull()} --- Envolvente convexa}
\begin{openscad}
hull() {
    sphere(r = 10);
    translate([30, 0, 0])
        sphere(r = 8);
}
\end{openscad}

\section{\texttt{minkowski()} --- Suma de Minkowski}
\begin{openscad}
// Redondear bordes de un cubo con una esfera
minkowski() {
    cube([10, 10, 1]);
    cylinder(r=2, h=1);
}
\end{openscad}

\section{\texttt{linear\_extrude()} --- Extrusión de 2D a 3D}
\begin{openscad}
linear_extrude(height = 10) { circle(r = 5); }
linear_extrude(height = 20, twist = 90) { square([10, 10]); }
\end{openscad}

\section{\texttt{text()} --- Texto 2D}
\begin{openscad}
text("SpiderBot", size = 10);
linear_extrude(height = 2)
    text("USS SpiderBot", size = 8, halign = "center");
\end{openscad}

\section{\texttt{offset()} --- Desplazamiento 2D}
\begin{openscad}
// Expandir con esquinas redondeadas
offset(r = 3) square([20, 20], center = true);
// Contraer y expandir para filete interno
offset(r = -2) offset(delta = 2)
    square([20, 20], center = true);
\end{openscad}

\section{\texttt{projection()} --- Proyección 3D a 2D}
\begin{openscad}
projection(cut = false) sphere(r = 10);  // Sombra
projection(cut = true)  sphere(r = 10);  // Corte transversal
\end{openscad}

\section{\texttt{polygon()} --- Perfiles 2D}
\begin{openscad}
polygon(points = [[0, 0], [10, 0], [5, 10]]);
linear_extrude(height = 3) {
    polygon(points = [[0, 0], [20, 0], [0, 10]]);
}
\end{openscad}

\section{\texttt{rotate\_extrude()} --- Revolución}
\begin{openscad}
rotate_extrude(angle = 360, convexity = 10) {
    translate([20, 0, 0]) circle(r = 5);
}
\end{openscad}

\section{Caso Real: El fémur con \texttt{hull()}}
\begin{openscad}
union() {
    hull() {
        cylinder(r = 9, h = 6, center=true);
        translate([link_len/2, -3, 1])
            cube([12, 12, 8], center=true);
    }
    translate([link_len - 12, -8, 0])
        cube([30, 28, link_t], center=true);
}
\end{openscad}

Gussets triangulares con \texttt{polygon()} + \texttt{linear\_extrude()}:
\begin{openscad}
translate([-22.25, cuerpo_r - 5.0, 5.5]) {
    rotate([0, -90, 0])
        linear_extrude(height = 2.0)
            polygon(points = [[0, 0], [11, 0], [0, 5]]);
}
\end{openscad}

% ═════════════════════════════════════════════════════════════════════════════
% CAPÍTULO 8
% ═════════════════════════════════════════════════════════════════════════════
\chapter{Archivos Múltiples --- \texttt{use} e \texttt{include}}

\section{\texttt{use <archivo.scad>}} Importa solo los módulos:
\begin{openscad}
// main.scad
use <tornillo.scad>
tornillo(20);  // El módulo está disponible
\end{openscad}

\section{\texttt{import()} --- Importar STL/SVG/DXF}
\begin{openscad}
import("pieza_externa.stl");
linear_extrude(height = 3) import("logo.svg");
import("perfil.dxf");
\end{openscad}

\section{\texttt{include <archivo.scad>}} Importa todo (menos usado):
\begin{openscad}
include <tornillo.scad>
// Todo el código de tornillo.scad se ejecuta aquí
\end{openscad}

\section{Organización de proyectos}
\begin{verbatim}
robot/
├── cuerpo.scad         // Chasis y soportes
├── femur.scad          // Eslabón de pata
├── tibia.scad          // Pierna inferior
└── ensamble.scad       // Ensamble completo (solo use)
\end{verbatim}

\section{Caso Real: El ensamble del SpiderBot}
\begin{openscad}
// cad/ensamble_cuadruepodo.scad
use <cuerpo_cuadruepdo.scad>
use <eslabon_pata_cuadrupedo.scad>
use <tibia_pata.scad>
$fn = 20;  // Prevalece para render más rápido

// Archivos "instanciadores":
// placa_base_inferior.scad:
use <cuerpo_cuadruepdo.scad>
placa_base_cuadruepodo();
\end{openscad}

% ═════════════════════════════════════════════════════════════════════════════
% CAPÍTULO 9
% ═════════════════════════════════════════════════════════════════════════════
\chapter{Renderizado, Debug y Exportación a STL}

\section{Preview (F5) vs Render (F6)}
\begin{center}
\begin{tabular}{lll}
\toprule
Aspecto & Preview (F5) & Render (F6) \\
\midrule
Motor & OpenCSG (GPU) & CGAL (CPU) \\
Velocidad & Instantáneo & Puede tomar minutos \\
Precisión & Aproximada & Exacta (manifold) \\
Exportable STL & No & Sí \\
\bottomrule
\end{tabular}
\end{center}

\section{Modifier Characters --- Debug visual}
\begin{center}
\begin{tabular}{cll}
\toprule
Carácter & Nombre & Efecto \\
\midrule
\texttt{\%} & Background & Gris translúcido, excluido de booleanas \\
\texttt{\#} & Debug & Rosa translúcido, incluido \\
\texttt{!} & Root & Solo este sub-árbol \\
\texttt{*} & Disable & Ignorado completamente \\
\bottomrule
\end{tabular}
\end{center}

\begin{openscad}
difference() {
    cube([30, 30, 10], center=true);
    #translate([0, 0, -1])          // Rosa = debug
        cylinder(r = 5, h = 12);
    %translate([10, 0, -1])         // Gris = background
        cylinder(r = 3, h = 12);
    *translate([-10, 0, -1])        // Ignorado
        cylinder(r = 3, h = 12);
}
\end{openscad}

\section{\texttt{render()} --- Forzar CGAL en preview}
\begin{openscad}
render(convexity = 4)
    difference() {
        cube([20, 20, 10], center=true);
        sphere(r = 8);
    }
\end{openscad}

\section{\texttt{\$t} y animación}
Vista > Animación > FPS=10, Steps=100:
\begin{openscad}
translate([0, 0, 10 * sin($t * 360)]) sphere(r = 5);
rotate([0, 0, $t * 360]) cube([20, 5, 5], center=true);
\end{openscad}

\section{\texttt{\$preview} --- Diferenciar preview de render}
\begin{openscad}
$fn = $preview ? 12 : 72;
cylinder(r = 10, h = 20);
\end{openscad}

\section{Exportar a STL (F7)}
Con F6 renderizado: Archivo > Exportar > Exportar como STL, o F7.

\section{CLI de OpenSCAD}
\begin{verbatim}
# Exportar pieza individual
openscad -o placa_base.stl cad/placa_base_inferior.scad
# Con parámetros personalizados
openscad -o femur_largo.stl -D "link_len=70" cad/eslabon_pata_cuadrupedo.scad
# Calidad específica
openscad -o pieza.stl -D "\$fn=100" archivo.scad
\end{verbatim}

\section{Parámetros de calidad}
\begin{center}
\begin{tabular}{lcc}
\toprule
Parámetro & Por defecto & Recomendado STL \\
\midrule
\texttt{\$fn} & (automático) & 50--80 \\
\texttt{\$fa} & 12° & 6°--2° \\
\texttt{\$fs} & 2mm & 0.5--1mm \\
\bottomrule
\end{tabular}
\end{center}

\section{Caso Real: Comandos CLI para el SpiderBot}
\begin{verbatim}
openscad -o placa_base_inferior.stl -D "\$fn=60" cad/placa_base_inferior.scad
openscad -o femur.stl            -D "\$fn=50" cad/eslabon_pata_cuadrupedo.scad
openscad -o tibia.stl            -D "\$fn=50" cad/tibia_pata.scad
openscad -o placa_superior.stl   -D "\$fn=60" cad/placa_base_superior.scad
openscad -o spiderbot_ensamble.stl -D "\$fn=30" cad/ensamble_cuadruepodo.scad
\end{verbatim}

% ═════════════════════════════════════════════════════════════════════════════
% CAPÍTULO 10
% ═════════════════════════════════════════════════════════════════════════════
\chapter{Proyecto Final --- Replica el USS SpiderBot desde Cero}

Ejercicio guiado. Las soluciones completas están en \texttt{cad/}.

\section{Paso 1: El chasis circular}
Crea \texttt{mi\_spiderbot.scad}:
\begin{enumerate}
    \item Define \texttt{\$fn = 50}.
    \item Variables: \texttt{cuerpo\_r = 65.0}, \texttt{cuerpo\_t = 3.0}, \texttt{spacer\_r = 55.0}.
    \item Módulo \texttt{placa\_base()} con \texttt{cylinder(r = cuerpo\_r, h = cuerpo\_t)}.
    \item Con \texttt{difference()}, agrega 4 agujeros en 0°, 90°, 180°, 270° a radio \texttt{spacer\_r}.
    \item Agrega 4 ranuras tipo velcro de \texttt{[4, 15, cuerpo\_t+2]} en (\texttt{±20, ±10}).
\end{enumerate}

\section{Paso 2: El soporte de servo paramétrico}
Módulo \texttt{soporte\_servo()}:
\begin{enumerate}
    \item Base: \texttt{cube([36, 28, 15], center=true)} en Z=7.5.
    \item Cavidad dual-depth: inferior a Z=2.9, superior a Z=9.5.
    \item Ranura para bridas en Y=5.5.
    \item Agujeros M2 con \texttt{rotate([90,0,0]) cylinder} en X=±14.25.
    \item Abertura delantera para corona en (5.5, 12.0, 7.5).
\end{enumerate}

\section{Paso 3: Integrar soportes en el chasis}
Dentro de \texttt{placa\_base()}, agrega:
\begin{openscad}
for (a = [45, 135, 225, 315]) {
    rotate([0, 0, a])
        translate([cuerpo_r - 7, 0, cuerpo_t/2])
            rotate([0, 0, -90])
                soporte_servo();
}
\end{openscad}

\section{Paso 4: El fémur con hull()}
Crea \texttt{femur.scad}. Define \texttt{link\_len}, \texttt{link\_h}, \texttt{link\_t}.
Usa \texttt{hull()} para conectar cilindro de cadera con bloque de rodilla.

\section{Paso 5: La tibia}
Crea \texttt{tibia.scad}. Define \texttt{tibia\_h}, \texttt{tibia\_t}.
Usa tres \texttt{hull()} consecutivos. Agrega \texttt{sphere(r=5.5)} para el pie.

\section{Paso 6: Ensamblar todo}
Crea \texttt{ensamble\_final.scad}:
\begin{enumerate}
    \item \texttt{use} para importar tus 3 archivos.
    \item Crea \texttt{dummy\_servo()} y \texttt{pata\_articulada()}.
    \item \texttt{ensamble\_completo()} con placa base, 4 patas y placa superior.
\end{enumerate}

\section{Paso 7: Exportar a STL}
\begin{verbatim}
openscad -o mi_placa_base.stl -D "\$fn=60" mi_spiderbot.scad
openscad -o mi_femur.stl      -D "\$fn=50" femur.scad
openscad -o mi_tibia.stl      -D "\$fn=50" tibia.scad
\end{verbatim}

% ═════════════════════════════════════════════════════════════════════════════
% APÉNDICE A
% ═════════════════════════════════════════════════════════════════════════════
\chapter*{Apéndice A: Hoja de Referencia Rápida}
\addcontentsline{toc}{chapter}{Apéndice A: Hoja de Referencia Rápida}

\begin{center}
\begin{longtable}{lll}
\toprule
Comando/Función & Sintaxis & Descripción \\
\midrule
\endhead
\bottomrule
\endfoot
\texttt{cube} & \texttt{cube([w,d,h], center)} & Cubo o prisma rectangular \\
\texttt{sphere} & \texttt{sphere(r)} & Esfera \\
\texttt{cylinder} & \texttt{cylinder(r,h)} / \texttt{(r1,r2,h)} & Cilindro o cono \\
\texttt{polyhedron} & \texttt{polyhedron(points,faces)} & Poliedro arbitrario \\
\texttt{translate} & \texttt{translate([x,y,z])} & Desplazar \\
\texttt{rotate} & \texttt{rotate([x,y,z])} & Rotar (grados) \\
\texttt{scale} & \texttt{scale([x,y,z])} & Escalar \\
\texttt{mirror} & \texttt{mirror([x,y,z])} & Reflejo especular \\
\texttt{resize} & \texttt{resize([x,y,z])} & Redimensionar exacto \\
\texttt{color} & \texttt{color("nombre")} & Color (solo preview) \\
\texttt{union} & \texttt{union() \{ \}} & Fusionar sólidos \\
\texttt{difference} & \texttt{difference() \{ A; B; \}} & Restar B de A \\
\texttt{intersection} & \texttt{intersection() \{ \}} & Intersección \\
\texttt{hull} & \texttt{hull() \{ \}} & Envolvente convexa \\
\texttt{minkowski} & \texttt{minkowski() \{ \}} & Suma Minkowski \\
\texttt{linear\_extrude} & \texttt{linear\_extrude(h)} & Extruir 2D a 3D \\
\texttt{rotate\_extrude} & \texttt{rotate\_extrude()} & Revolución 2D \\
\texttt{text} & \texttt{text("str", size)} & Texto 2D \\
\texttt{polygon} & \texttt{polygon(points=[...])} & Polígono 2D \\
\texttt{offset} & \texttt{offset(r)} / \texttt{delta} & Desplazar contorno 2D \\
\texttt{projection} & \texttt{projection(cut)} & Proyectar 3D a 2D \\
\texttt{for} & \texttt{for(i=[inicio:fin])} & Bucle (unión) \\
\texttt{intersection\_for} & \texttt{intersection\_for(...)} & Bucle (intersección) \\
\texttt{if} & \texttt{if(cond)\{\}else\{\}} & Condicional \\
\texttt{let} & \texttt{let(v=val)\{\}} & Variables locales \\
\texttt{assert} & \texttt{assert(cond,"msg")} & Validación \\
\texttt{module} & \texttt{module nom(p)\{\}} & Definir módulo \\
\texttt{function} & \texttt{function nom(p)=expr} & Función (devuelve valor) \\
\texttt{children} & \texttt{children(i)} & Hijos del módulo \\
\texttt{use} & \texttt{use <archivo.scad>} & Importar módulos \\
\texttt{include} & \texttt{include <archivo.scad>} & Importar todo \\
\texttt{import} & \texttt{import("arch.stl")} & Importar STL/SVG/DXF \\
\texttt{echo} & \texttt{echo(valor)} & Imprimir en consola \\
\texttt{str} & \texttt{str(a,b,...)} & Concatenar a string \\
\texttt{concat} & \texttt{concat(a,b,...)} & Concatenar listas \\
\texttt{len} & \texttt{len(lista)} & Longitud \\
\texttt{[for]} & \texttt{[for(i=[0:N]) expr]} & List comprehension \\
\texttt{\$fn} & \texttt{\$fn = N} & Fragmentos resolución \\
\texttt{\$fa} & \texttt{\$fa = ang} & Ángulo mínimo \\
\texttt{\$fs} & \texttt{\$fs = tam} & Tamaño mínimo \\
\texttt{\$t} & \texttt{\$t} (lectura) & Animación (0 a 1) \\
\texttt{\$preview} & \texttt{\$preview} & true en F5, false en F6 \\
\texttt{render} & \texttt{render(convexity=N)} & Forzar malla CGAL \\
\texttt{\%} & \texttt{\% \{ \}} & Background (gris) \\
\texttt{\#} & \texttt{\# \{ \}} & Debug (rosa) \\
\texttt{!} & \texttt{! \{ \}} & Root (solo esto) \\
\texttt{*} & \texttt{* \{ \}} & Disable (ignorar) \\
\end{longtable}
\end{center}

% ═════════════════════════════════════════════════════════════════════════════
% APÉNDICE B
% ═════════════════════════════════════════════════════════════════════════════
\chapter*{Apéndice B: Recursos y Documentación Oficial}
\addcontentsline{toc}{chapter}{Apéndice B: Recursos y Documentación Oficial}

\begin{itemize}
    \item \textbf{Documentación oficial de OpenSCAD:} \url{https://openscad.org/documentation.html}
    \item \textbf{Cheat Sheet oficial:} \url{https://openscad.org/cheatsheet/}
    \item \textbf{Manual de usuario (Wikibooks):} \url{https://en.wikibooks.org/wiki/OpenSCAD_User_Manual}
    \item \textbf{Reddit:} \url{https://reddit.com/r/openscad}
    \item \textbf{GitHub:} \url{https://github.com/openscad/openscad}
    \item \textbf{Librería MCAD:} \url{https://github.com/openscad/MCAD}
    \item \textbf{dotSCAD (utilidades):} \url{https://github.com/JustinSDK/dotSCAD}
    \item \textbf{Repositorio USS SpiderBot:} \url{https://github.com/moises-inc/USS-SPIDERBOT-solemne-3}
\end{itemize}

\vspace{2cm}

\begin{center}
\rule{5cm}{0.4pt}

\small\textit{Guía generada para el proyecto USS SpiderBot --- Taller de Programación I, Universidad San Sebastián. 2026.}

\small\textit{Los archivos de referencia están en \texttt{cad/} del repositorio.}
\end{center}

\end{document}
